\chapter{Préparer sa machine}

\section{Ce qu'il vous faut vraiment}

\begin{center}
\begin{tabular}{p{4.2cm}p{9.8cm}}
\toprule
\textbf{Élément} & \textbf{Exigence réelle} \\
\midrule
Carte graphique & 8~Go de mémoire vidéo suffisent pour un modèle de 7~milliards
                  de paramètres quantifié. C'est la configuration sur laquelle
                  toutes les mesures de ce cours ont été prises. \\
Mémoire vive & 16~Go confortables. Le corpus complet (25~Mo de texte) ne pèse
               rien à côté du modèle. \\
Disque & 15~Go : le modèle (4,4~Go), les checkpoints (231~Mo chacun, on en
         garde deux), et de la marge. \\
Système & Windows, Linux ou macOS. \\
\bottomrule
\end{tabular}
\end{center}

\begin{nkpiege}
Le facteur limitant n'est pas la \emph{vitesse} de votre carte, mais sa
\textbf{mémoire}. Une carte rapide avec 4~Go ne chargera pas le modèle ; une
carte lente avec 8~Go le fera tourner. Vérifiez la mémoire avant tout le reste.
\end{nkpiege}

\section{Comment la machine parle à la carte graphique}

Les calculs d'entraînement ne se font pas sur le processeur : ils se font sur la
carte graphique, qui sait faire des millions de multiplications à la fois. Pour
lui parler, il faut une \textbf{interface} — un langage commun. Nkentseu en
gère quatre.

\begin{center}
\begin{tabular}{lp{9.4cm}}
\toprule
\textbf{Interface} & \textbf{Où elle existe} \\
\midrule
Vulkan  & Windows, Linux, macOS (via MoltenVK). \textbf{C'est celle qu'on
          utilise.} \\
DirectX 11 & Windows uniquement. \\
DirectX 12 & Windows uniquement. \\
OpenGL  & Partout, mais la moins adaptée aux calculs. \\
\bottomrule
\end{tabular}
\end{center}

\subsection{Pourquoi Vulkan, et pas DirectX}

Pour une raison très concrète. Le modèle contient un tenseur particulier —
celui qui produit le mot final, appelé \texttt{output.weight} — qui pèse
\textbf{426~Mo à lui seul}. Or DirectX~11 ne garantit que 128~Mo par ressource.

Le programme s'en occupe tout seul : il se verrouille sur Vulkan au démarrage et
vous l'annonce.

\begin{nkcode}[]{Ce que vous verrez au lancement}
backend compute : Vulkan
  [NkQwen2LoraGpu] backend = Vulkan
  [NkQwen2LoraGpu] L=28 d=3584 ffn=18944 nH=28 nKV=4 hd=128 vocab=152064
\end{nkcode}

\begin{nkpiege}
Ne concluez jamais d'une mesure sans avoir lu la ligne « backend compute ». Une
même commande peut tourner sur deux interfaces différentes et donner des durées
sans rapport. Un résultat dont on ignore le backend est un résultat sans valeur.
\end{nkpiege}

\section{Installer les outils}

\subsection{Sur Windows}

Installez le \textbf{SDK Vulkan} depuis \url{https://vulkan.lunarg.com}. C'est
tout : le reste est fourni par le dépôt.

Vérifiez ensuite que la variable \texttt{VULKAN\_SDK} existe :

\begin{nkcode}[]{PowerShell}
echo $env:VULKAN_SDK
\end{nkcode}

Si la commande n'affiche rien, redémarrez la session : l'installateur pose la
variable, mais les fenêtres déjà ouvertes ne la voient pas.

\subsection{Sur Linux}

\begin{nkcode}[]{Debian, Ubuntu et dérivées}
sudo apt install libvulkan-dev vulkan-tools libx11-dev libxext-dev
vulkaninfo | head -20
\end{nkcode}

Si \texttt{vulkaninfo} affiche le nom de votre carte, vous êtes prêt.

\begin{nkpiege}
\textbf{Le piège du mode « headless ».} Sur un serveur sans écran, on est tenté
de compiler en \texttt{-{}-linux-backend=headless}. Ne le faites pas pour
l'entraînement : ce mode ne relie pas Vulkan, donc votre programme n'aurait
\emph{aucune} carte graphique à sa disposition. Compilez en \texttt{xlib} — le
mode par défaut — même si vous n'affichez aucune fenêtre.
\end{nkpiege}

\subsection{Sur macOS}

Vulkan y passe par MoltenVK, une couche qui traduit vers Metal. Installez le SDK
Vulkan pour macOS depuis le même site que pour Windows.

\subsection{Sur WSL2 (Windows Subsystem for Linux)}

Cela fonctionne, mais lisez ceci avant de vous en réjouir : WSL2 utilise
\emph{la même carte graphique} que Windows, avec une couche de traduction en
plus. Vous serez donc \textbf{plus lent} qu'en Windows natif, jamais plus
rapide.

WSL2 sert à vérifier que votre travail fonctionne sous Linux sans acheter une
seconde machine. Pour un entraînement de plusieurs jours, restez sur Windows
natif ou sur une vraie machine Linux.

\section{Compiler les outils}

Tout se compile avec \texttt{jenga}, depuis la racine du dépôt.

\begin{nkcode}[]{Les deux programmes de ce cours}
jenga build --target NKQwen2Train --config Release
jenga build --target NKQwen2Chat  --config Release
\end{nkcode}

\texttt{Release} et non \texttt{Debug} : la version de mise au point est
plusieurs fois plus lente, ce qui sur dix jours d'entraînement fait une
différence de plusieurs semaines.

Vous devez lire, à la fin :

\begin{nkcode}[]{Compilation réussie}
Projects Built:  21/21
Status:          SUCCESS
\end{nkcode}

\begin{nkpiege}
Si vous compilez \emph{tout} le dépôt avec \texttt{jenga build} sans
\texttt{-{}-target}, vous verrez peut-être échouer un projet de démonstration
sans rapport avec l'IA. Ce n'est pas votre faute et cela n'empêche rien :
compilez toujours la cible qui vous intéresse.
\end{nkpiege}

\section{Se procurer un modèle de départ}

Il vous faut un modèle déjà entraîné, au format \textbf{GGUF}. Le plus simple
est de passer par Ollama :

\begin{nkcode}[]{Télécharger Qwen2.5 7B}
ollama pull qwen2.5:7b
\end{nkcode}

Le fichier atterrit dans un dossier de blobs, sous un nom qui est son empreinte :

\begin{center}
\begin{tabular}{lp{8.6cm}}
\toprule
Windows & \texttt{\%USERPROFILE\%\textbackslash.ollama\textbackslash models\textbackslash blobs\textbackslash} \\
Linux, macOS & \texttt{\textasciitilde/.ollama/models/blobs/} \\
\bottomrule
\end{tabular}
\end{center}

Repérez le fichier de 4,4~Go : c'est le modèle. Les petits fichiers à côté sont
des métadonnées.

\begin{nkpiege}
\textbf{Vérifiez que vous avez le bon fichier.} Un dossier de blobs contient
souvent plusieurs modèles de taille voisine. Se tromper de blob donne un modèle
qui répond n'importe quoi, avec des jetons spéciaux inattendus — et l'on cherche
l'erreur pendant des heures dans son code alors qu'elle est dans le chemin.
Notez le chemin exact de \emph{votre} modèle, une bonne fois.
\end{nkpiege}

\section{Vérifier que tout fonctionne}

Avant d'entraîner quoi que ce soit, faites parler le modèle. Si cette étape
échoue, inutile d'aller plus loin.

\begin{nkcode}[]{Windows — PowerShell}
.\Build\Bin\Release-Windows\NKQwen2Chat\NKQwen2Chat.exe `
  "$env:USERPROFILE\.ollama\models\blobs\sha256-VOTRE_BLOB"
\end{nkcode}

\begin{nkcode}[]{Linux et macOS}
./Build/Bin/Release-Linux/NKQwen2Chat/NKQwen2Chat \
  ~/.ollama/models/blobs/sha256-VOTRE_BLOB
\end{nkcode}

Le modèle met une dizaine de secondes à se charger, puis vous rend la main.
Posez une question, appuyez sur Entrée, et regardez la réponse s'écrire mot à
mot.

\begin{nkretenir}
Trois vérifications, dans cet ordre, avant tout entraînement : la ligne
« backend compute » dit \textbf{Vulkan} ; le modèle se charge sans erreur de
mémoire ; il répond quelque chose de sensé. Si les trois passent, votre machine
est prête.
\end{nkretenir}

\begin{nkexo}
Chargez le modèle et posez-lui trois questions : une factuelle (« Quelle est la
capitale du Cameroun ? »), une de raisonnement (« Si tous les A sont B et que
X est un A, que peut-on dire de X ? »), une créative (« Écris deux phrases sur
la pluie »). Notez les réponses dans un fichier. Vous les comparerez à celles
d'après l'entraînement, au chapitre~7 — et c'est la seule façon honnête de
juger si votre entraînement a servi à quelque chose.
\end{nkexo}
